\bumppoints{24}

\question
A harmless fly looks like a dangerous wasp. The fly is using

\begin{choices}
\choice camouflage. 
\choice disruptive coloration. 
\choice warning coloration. 
\correctchoice mimicry. 
\choice cryptic coloration. 
\end{choices}

\question
Which of the following statements is consistent with the principle of competitive exclusion?  

\begin{choices}
\choice Evolution tends to increase competition between related species.   
\correctchoice Even a slight reproductive advantage will eventually lead to the elimination of the less well adapted of two competing species. 
\choice The density of one competing species will have a positive affect on the population growth of the other competing species. 
\choice Bird species generally do not compete for nesting sites.  
\choice A keystone species excludes the dominant competitor from a community.   
\end{choices}

\question
While observing a population of lizards during one week, you notice that the number of adults is higher at the end of the week than at the start. The number of juveniles is unchanged. The most likely explanation for such an observation would be 

\begin{choices}
\choice decreased death rate. 
\choice decreased emigration.  
\choice increased birth rate.  
\correctchoice increased immigration.  
\choice decreased immigration. 
\choice increased emigration.  
\end{choices}

\question
Given the following values and assuming exponential population growth, what will be the total population size after one generation has elapsed? Round to nearest whole individual.\vspace{1ex}

\hspace{\leftmargin}%
\begin{minipage}{2in}
$b = 0.09$ \qquad $N = 18,657$\\
$d = 0.12$ \qquad $K = 15,000$\\
\end{minipage}

\begin{choices}
\choice $560$ %dontshuffle 
\choice $3,657$ %dontshuffle 
\choice $15,000$ %dontshuffle 
\correctchoice $18,097$ %dontshuffle 
\choice $19,217$ %dontshuffle 
\end{choices}

\newpage

\question
Which of these statements about human evolution is correct?  

\begin{choices}
\correctchoice Different species of the genus \textit{Homo} have coexisted at various times throughout hominin evolution.   
\choice The evolution of upright posture and enlarged brain occurred simultaneously.   
\choice The ancestors of \textit{Homo sapiens} were chimpanzees.   
\choice \textit{Australopithecus} was the first group to use projectile hunting tools. 
\choice Human evolution has proceeded in an linear fashion from \textit{Australopithecus} to \textit{Homo sapiens.}   
\end{choices}

\question
Which of the following is most likely to contribute to density-dependent regulation of populations?

\begin{choices}
\correctchoice Intraspecific competition for nutrients.
\choice The removal of toxic waste by decomposers.
\choice Earthquakes.
\choice Floods.
\choice Fires.
\end{choices}

\question
Which of the following is \emph{not} directly concerned with community ecology?

\begin{choices}
\correctchoice intraspecific competition
\choice resource partitioning
\choice symbiosis
\choice herbivory
\choice predation
\end{choices}

\question
A small population of white-footed mice has the same intrinsic rate of increase (\textit{r}) as a large population.  If everything else is equal, then

\begin{choices}
\choice the growth trajectories of the two populations will proceed in opposite directions. 
\choice the number of individuals added will depend on which population has the higher birth rate. 
\choice the small population will add more individuals per unit time. 
\correctchoice the large population will add more individuals per unit time.  
\choice the two populations will add equal numbers of individuals per unit time.
\end{choices}

\question
Several species of darters (a group of small fishes) live together on the gravel bottoms of small streams. They use similar resources in slightly different ways. The darters can live in the same habitat type  because they

\begin{choices}
	\choice have the same niches within the stream.
	\choice feed at different times of the day.
	\correctchoice partitioned the resources to reduce competition.
	\choice can swim at different heights above the gravel bottom.
	\choice out-competed other species of fishes, such as catfishes.
\end{choices}

\newpage

\question
The attribute that most influences population density and rates of immigration and emigration among habitat patches in a metapopulation are

\begin{choices}
\correctchoice patch size, patch isolation and habitat quality. 
\choice habitat quality, size of sink patches and isolation of source patches. 
\choice patch size, patch isolation and species diversity. 
\choice patch isolation, habitat quality, and size of source patches. 
\choice patch size, habitat quality, and isolation of sink patches. 
\end{choices}

\question
The dispersion pattern for individuals across the entire metapopulation shown in question \ref{ques:metapopulation} is most likely

\begin{choices}
\choice uniform. 
\choice even. 
\choice gene flow. 
\correctchoice clumped. 
\choice random. 
\end{choices}

\question
In the experiment shown in the figure below, \textit{Balanus} was removed from the habitat shown on the left. Which of the following statements is a valid conclusion of this experiment?\\
\includegraphics[width=0.9\textwidth]{exam3_balanus}

\begin{choices}
\choice \textit{Balanus} can survive only in the lower intertidal zone because it is unable to resist dehydration.
\choice These two species of barnacle do not show competitive exclusion.
\choice \textit{Balanus} is inferior to \textit{Chthamalus} in competing for space on rocks lower in the intertidal zone. 
\correctchoice The removal of \textit{Balanus} shows that the realized niche of \textit{Chthamalus} is smaller than its fundamental niche.
\choice If \textit{Chthamalus} were removed, \textit{Balanus}'s fundamental niche would become larger.
\end{choices}

\vfill

Hi!